\documentclass{article}
\usepackage{ctex}
\usepackage{amsmath,amssymb}
\usepackage{graphicx}
\usepackage{geometry}
\geometry{a4paper,margin=2.5cm}
\usepackage{hyperref}
\hypersetup{colorlinks=true,linkcolor=blue}

\title{基于波动力学的神经网络类别偏置预测模型研究}
\author{}
\date{}

\begin{document}
\maketitle

\section{绪论}
\subsection{研究背景与意义}
随着深度学习技术的快速发展，神经网络在图像分类、目标检测等任务中取得了优异效果。但现有模型普遍存在类别偏置问题：训练集中不同类别样本数量不均衡、特征分布差异大，导致模型对少数类识别精度偏低，泛化能力下降。

波动力学从波动、演化的角度描述状态传播规律，和神经网络特征逐层传递、信息演化的过程存在内在相似性。将波动力学思想引入神经网络，用来建模类别之间的特征传播、偏差累积过程，对预测、缓解类别偏置有重要研究价值。

本研究旨在构建结合波动力学思想的神经网络类别偏置预测模型，提前预判不同类别上的偏差程度，为后续的样本重加权、数据增强、损失函数优化提供依据，提升不均衡数据集下模型的稳定性与分类性能。

\subsection{国内外研究现状}
在类别不均衡与类别偏置方向，现有方案主要分为数据层面与算法层面：数据层面包含重采样、过采样、生成少数类样本；算法层面包括加权损失、迁移学习、解耦训练等。这类方法大多是事后修正，缺少对偏置发展趋势的预测机制。

波动力学相关方法在机器学习中多用于波动方程求解、图信号传播、神经波动网络等方向，多用于图学习、物理仿真领域，将其用于类别偏置预测的研究相对较少。如何结合波动演化规律刻画类别偏差随训练过程的变化，是一个值得探索的方向。

\subsection{研究内容与技术路线}
本文主要研究内容如下：
\begin{enumerate}
    \item 分析不均衡数据集下神经网络类别偏置产生机理，定义可量化的类别偏置指标；
    \item 借鉴波动力学波动演化思想，搭建类别偏置预测模型，建模偏差随训练迭代的传播过程；
    \item 在公开不均衡分类数据集上开展对比实验，验证模型对类别偏置的预测效果；
    \item 根据预测结果，设计对应的偏置抑制策略，对比基线模型，分析分类精度提升效果。
\end{enumerate}

技术路线：数据集构建与偏置指标定义 → 波动特征建模与预测网络设计 → 模型训练与预测结果验证 → 偏置抑制策略设计 → 实验对比与结果分析。

\subsection{论文结构安排}
第一章为绪论，介绍研究背景、意义、国内外现状、研究内容与论文框架；
第二章介绍相关理论基础，包括类别不均衡、类别偏置评价指标、波动力学基础、基础神经网络；
第三章提出基于波动力学的类别偏置预测模型，给出网络结构与损失设计；
第四章设计实验，介绍数据集、评价指标、对比方案，分析实验结果；
第五章总结全文工作，分析不足并展望后续研究方向。

\section{相关理论基础}
\subsection{类别不均衡与类别偏置}
类别不均衡指数据集中不同类别的样本数量差距较大，分为天然不均衡、人为采样不均衡等情况。类别不均衡往往会带来类别偏置问题：模型在训练时会偏向样本数量多的大类，在少数类别上决策边界发生偏移，预测输出存在系统性偏差。

类别偏置和单纯的数据不平衡有区别：即使样本数量均衡，若各类特征分布复杂度不同，依然会产生类别偏置。常用评价指标包括召回率、F1值、均衡准确率、偏差方差分解等。在不均衡分类任务中，不能只依靠总体准确率评价模型效果，需要重点关注各类别上的性能差异。

\subsection{神经网络特征传播机制}
深度神经网络由多层非线性变换堆叠而成，输入样本经过逐层卷积、激活、池化或者全连接运算，特征不断被提取、映射、抽象。这个过程可以看作信息在网络中的传播演化过程：浅层提取边缘、纹理等底层特征，深层学习语义、类别相关的高层特征。

在训练迭代过程中，梯度反向传播不断更新网络参数，各类特征的分布持续变化。类别偏置也会随着迭代逐步累积：大类样本的梯度主导参数更新，不断挤压少数类的特征空间，最终造成偏置不断放大。

\subsection{波动力学基础}
波动力学以波动方程为核心，描述物理量随时间、空间的传播与演化规律。一维波动方程形式为：
\[
\frac{\partial^2 u}{\partial t^2}=c^2\frac{\partial^2 u}{\partial x^2}
\]
其中 \(u(x,t)\) 为波动状态量，\(c\) 为波速。波动具有传播、叠加、衰减、色散等特性，状态会随时间逐步扩散、演化。

本研究借鉴波动演化思想，不直接求解物理波动方程，而是将类别偏差视作一种“波动信号”：把训练迭代当作时间维度，类别之间的特征关联当作空间连接，刻画偏差在训练过程中传播、累积、衰减的规律，以此实现对类别偏置变化趋势的预测。

\subsection{不均衡学习常用方法}
1. 数据层面：重采样，包括对少数类过采样、对多数类欠采样；生成式方法利用GAN、扩散模型合成少数样本。
2. 损失层面：加权交叉熵、Focal Loss、LDAM等，通过调整各类损失权重抑制偏置。
3. 训练策略层面：解耦训练、两阶段训练、迁移学习，先学习通用特征再微调少数类。

上述方法大多是补偿策略，缺少对偏置演化趋势的预测能力。本文正是基于这一点，引入波动演化模型提前预测偏置，为动态补偿提供依据。

\end{document}
